\documentclass[10pt,a4paper]{amsart}
\usepackage[margin=19mm]{geometry}
\usepackage{amsmath,amssymb,mathtools}
\usepackage[scr=rsfs,cal=euler]{mathalfa}
\usepackage{xfrac}
\usepackage[unicode,hidelinks,bookmarks=false]{hyperref}
\newcommand{\ie}{i.e.\ }
\newcommand{\cf}{cf.\ }
\newcommand{\resp}{respectively\ }
\newcommand{\Subsection}{\subsection}
\providecommand{\nobreakdash}{\nobreak\hbox{-}}
\setlength{\emergencystretch}{2em}
\multlinegap0pt
\usepackage{amssymb}
\usepackage{tikz}
\usepackage{mathtools}
\newtheorem{theorem}{Theorem}
\usepackage{cleveref}
\crefname{theorem}{Theorem}{Theorems}
\newcommand{\R}{\mathbb{R}}
\title{Information Geometry on the $\ell^2$-Simplex via the $q$-Root Transform}
\author{Levin Maier}
\date{Source: arXiv:2506.00485v2 (CC BY 4.0)}
\begin{document}
\maketitle

\noindent\emph{Source and licence.} Information Geometry on the $\ell^2$-Simplex via the $q$-Root Transform. Version arXiv:2506.00485v2 (CC BY 4.0), distributed by arXiv under the Creative Commons Attribution 4.0 International licence, http://creativecommons.org/licenses/by/4.0/, as recorded in the arXiv licence metadata for this exact version and shown on the abstract page https://arxiv.org/abs/2506.00485v2. Full licence notice: \texttt{licenses/arxiv-2506.00485-CC-BY-4.0.txt}.\par\smallskip

\noindent\textit{Editorial scope.} Counted unit: the article's theorem system (source0.tex lines 291--294). The article's complete original proof of it is delivered with it (source0.tex lines 295--311, one \texttt{proof} environment of 17 lines). No mathematics is written, repaired or re-derived: the statement and the proof are the article's own lines, in the article's own notation, and no retained line is substituted or re-flowed.  Retained uncounted as the source-local context the proof needs: lines 147--153; lines 214--216; lines 369--369.  Nothing was omitted from any retained span, so the package carries no omission record.  No retained line contains a citation, so the wrapper carries no bibliography and no bibliography entry.  No retained line contains a figure environment, an image-inclusion command or drawing code, so no image is delivered and the package declares no asset dependency.  Editorial formatting only, in addition: an AMS \texttt{amsart} preamble that restates the article's own declarations and macros used by the retained lines, namely the environments \texttt{theorem} on separate counters, together with the standard packages those lines need.  The retained environments print in this document as Theorem 1 in order of appearance, whereas the article prints the same items with the numbering of its own full document.  The complete native source of the article as distributed in the arXiv e-print for this exact version is bundled unchanged as \texttt{sources/179085/source0.tex}.
\begin{theorem}\label{t: square root map is isometry}
The square-root map \( \varPhi \) defined by
\[
\varPhi: (\Delta, \mathcal{G}^{\mathrm{FR}}) \longrightarrow (\mathcal{U}, \mathcal{G}^{\mathrm{L}^2}), \quad (p_n) \mapsto (\sqrt{p_n}),
\]
is an isometry.
\end{theorem}

\begin{equation}\tag{LP}\label{eq:ell2_LPP}
	\max_{(p_n)\in \bar{\Delta}} \langle (c_n), (p_n) \rangle_{\ell^2_\R}, \quad (c_n)\in \ell_2^\R.
\end{equation}

\section{Proof of \Cref{t: integrability of Hamiltonian system}}\label{Appendix proof of integrability of Hamiltonian system}

\begin{theorem}\label{t: integrability of Hamiltonian system}
The Hamiltonian system \( (\mathbb{CP}^\infty, \Omega^{\mathrm{FS}}, H_{(c_n)}) \) is completely integrable; that is, it possesses infinitely many Poisson-commuting conserved quantities. \\
Moreover let \( (x_n)(t) \) be a gradient flow line of the restriction of the gradient flow of \( H_{(c_n)} \) on \( (\mathbb{CP}^\infty, \mathcal{G}^{\mathrm{FS}}) \) to \( \mathcal{U} \). Then the limit $p_{\max} := \lim_{t \to \infty} \varPhi\left((x_n)(t)\right)$ exists and solves \eqref{eq:ell2_LPP}.
\end{theorem}

\begin{proof}
 For each \( n \in \mathbb{N} \), we define the sequence $(b_m)$ by 
 \[
b_m:=\begin{cases}
    c_n\, &, \quad\text{for} \quad m=n\\
    0\, &, \quad \text{else}
\end{cases}
 \]
and by this the Hamiltonian
\begin{equation}\label{eq: Hamiltonians on CP inf for integrability}
	H_n : \mathbb{C}P^\infty \longrightarrow \mathbb{R}, \quad [z_m] \mapsto  \sum_{n=0}^{\infty}b_m |z_m|^2= c_n|z_n|^2\, .
\end{equation}
From which we observe that the Hamiltonian flow \( \varphi_{H_n} \) fixes all coordinates except the \( n \)-th, where it acts as a phase rotation.  
Moreover, it is a straightforward exercise in symplectic geometry to verify that all the Hamiltonians \( H_n \) Poisson commute—both with each other and with \( H_{(c_n)} \). The precise details of this construction are carried out in \Cref{Appendix proof of integrability of Hamiltonian system}.\\
By identifying \( \mathcal{U} \) with its image in \( \mathbb{CP}^\infty \), and using \Cref{t: square root map is isometry}, along with the identity $H_{(c_n)} \circ \varPhi = F_{(c_n)}$ on $\Delta$,
where \( (F_{(c_n)}) \) is defined as in the line following \eqref{eq:ell2_LPP}, the proof is complete.
\end{proof}

\end{document}
